\chapter{Verhandlung und Koalitionen}

\section{Wenn Absprachen erlaubt sind}

\begin{merksatz}
\boxtitel{Der Leitfall: eine Budgetverhandlung}
Klinikum und Marienhospital teilen einen gemeinsam erwirtschafteten Betrag von $100$ Einheiten. Kommt keine Einigung zustande, so erhält das Klinikum $20$ und das Marienhospital $30$ aus Einzelverträgen.

\begin{center}
{\small
\renewcommand{\arraystretch}{1.5}
\begin{tabular}{@{}P{2.4cm}P{2.2cm}P{3.0cm}@{}}
\toprule
\textbf{Größe} & \textbf{Wert} & \\
\midrule
Gesamtbetrag & $100$ & \\
Konfliktpunkt & $(20,\ 30)$ & ohne Einigung\\
Überschuss & $50$ & zu verteilen\\
\bottomrule
\end{tabular}}
\end{center}

\begin{center}
\begin{minipage}{0.9\textwidth}
\textbf{Die Frage ist neu:} Nicht, was geschieht, sondern was eine \kursiv{angemessene} Aufteilung wäre.
\end{minipage}
\end{center}

\begin{center}
\begin{minipage}{0.88\textwidth}
\kursiv{Bisher beschrieb dieser Band, wo eigennützige Spieler landen. \textbf{Die kooperative Spieltheorie fragt anders}: Man gibt Forderungen an eine Lösung vor und prüft, welche sie erfüllt. Von Neumanns Einwand gegen Nashs Zugang --- Kapitel~4 --- zielte genau auf diesen Unterschied, und beide Zweige bestehen seitdem nebeneinander.}
\end{minipage}
\end{center}
\end{merksatz}

\section{Die Verhandlungslösung}

\begin{merksatz}
\boxtitel{Die vier Forderungen}
Gegeben sei eine Menge $S$ erreichbarer Nutzenpaare und ein Konfliktpunkt $\vek d$.

\begin{center}
{\small
\renewcommand{\arraystretch}{1.55}
\begin{tabular}{@{}P{2.6cm}P{6.0cm}@{}}
\toprule
\textbf{Pareto} & Es gibt keinen Punkt in $S$, der beiden mehr bietet.\\
\textbf{Symmetrie} & Sind die Lage und die Ausgangspunkte für beide gleich, so ist es das Ergebnis auch.\\
\textbf{Skaleninvarianz} & Rechnet man einen Nutzen affin um, so rechnet sich das Ergebnis mit um --- Satz 4.1.\\
\textbf{Unabhängigkeit} & Entfernt man aus $S$ Punkte, die nicht gewählt wurden, so ändert sich das Ergebnis nicht.\\
\bottomrule
\end{tabular}}
\end{center}
\end{merksatz}

\begin{merksatz}
\boxtitel{Satz 14.1 --- Nash, 1950}
Es gibt genau eine Lösung, die alle vier Forderungen erfüllt:
\[ \vek u^{*} = \operatorname*{arg\,max}_{\vek u \in S,\ \vek u \ge \vek d} \ (u_1-d_1)\,(u_2-d_2) . \]

\begin{center}
\begin{minipage}{0.9\textwidth}
\textbf{Man maximiert das Produkt der Zugewinne} --- nicht deren Summe.
\end{minipage}
\end{center}

\schritt{Der Leitfall}
Mit $u_1+u_2 = 100$ und $\vek d = (20,30)$:
\[ f(u_1) = (u_1-20)\,\bigl(100-u_1-30\bigr) = (u_1-20)(70-u_1) \]
\[ f'(u_1) = (70-u_1) - (u_1-20) = 90-2u_1 = 0 \quad\Longrightarrow\quad u_1 = 45 \]

\[ \vek u^{*} = (45,\ 55), \qquad \text{Produkt } 25\cdot25 = 625 \]

\begin{center}
\begin{minipage}{0.9\textwidth}
\textbf{Der Überschuss von $50$ wird hälftig geteilt} --- jeder erhält seinen Konfliktwert plus $25$.
\end{minipage}
\end{center}

\begin{center}
\begin{minipage}{0.88\textwidth}
\kursiv{Bei linearer Grenze läuft die Lösung stets darauf hinaus. \textbf{Die Ungleichheit des Ergebnisses ($45$ gegen $55$) stammt vollständig aus den Konfliktpunkten}, nicht aus der Verhandlung. Wer seine Ausgangslage verbessert, verbessert das Ergebnis eins zu eins --- und das ist der Grund, warum in Verhandlungen so viel Aufwand darauf verwendet wird, die eigene Rückfallmöglichkeit darzustellen.}
\end{minipage}
\end{center}

\schritt{Warum das Produkt}
\begin{center}
\begin{minipage}{0.88\textwidth}
\kursiv{Die Summe der Zugewinne wäre nicht skaleninvariant: Rechnet man den Nutzen des einen in andere Einheiten um, so verschöbe sich das Maximum. \textbf{Das Produkt ist invariant}, weil ein Faktor $\alpha$ es nur mit $\alpha$ multipliziert und die Stelle des Maximums unverändert lässt. Die dritte Forderung erzwingt also die Gestalt der Lösung fast allein.}
\end{minipage}
\end{center}

\schritt{Die strittige Forderung}
\begin{center}
\begin{minipage}{0.88\textwidth}
\kursiv{Es ist wieder die Unabhängigkeit --- wie bei Arrow in Kapitel~12. \textbf{Gibt man sie auf und verlangt stattdessen, dass beide denselben Anteil ihres jeweils Erreichbaren bekommen}, so erhält man die Lösung von Kalai und Smorodinsky (1975), die anders ausfällt. Es gibt keinen zwingenden Grund für die eine oder die andere Forderung.}
\end{minipage}
\end{center}
\end{merksatz}

\section{Verhandeln als Spielbaum}

\begin{merksatz}
\boxtitel{Das Modell von Rubinstein}
Beide machen abwechselnd Angebote. Jede Ablehnung kostet eine Runde; der Wert schrumpft mit $\delta_1$ bzw.\ $\delta_2$ je Runde.

\schritt{Satz 14.2 --- Rubinstein, 1982}
Es gibt genau ein teilspielperfektes Gleichgewicht. Wer das erste Angebot macht, erhält
\[ x^{*} = \frac{1-\delta_2}{1-\delta_1\delta_2} . \]

\schritt{Herleitung}
Sei $x$ der Anteil, den Spieler 1 als Anbieter behält, und $y$ der, den Spieler 2 als Anbieter behält.

\kursiv{Spieler 1 muss Spieler 2 gerade so viel bieten, dass dieser nicht ablehnt.} Lehnt Spieler 2 ab, so erhält er in der nächsten Runde als Anbieter den Wert $\delta_2\,y$:
\[ 1 - x = \delta_2\,y \]

\kursiv{Ebenso umgekehrt:}
\[ 1 - y = \delta_1\,x \]

\kursiv{Einsetzen:}
\[ x = 1 - \delta_2\,(1-\delta_1 x) = 1-\delta_2+\delta_1\delta_2\,x \]
\[ x\,(1-\delta_1\delta_2) = 1-\delta_2 \quad\Longrightarrow\quad x = \frac{1-\delta_2}{1-\delta_1\delta_2} \]
$\blacksquare$

\schritt{Die Werte}
\begin{center}
{\small
\renewcommand{\arraystretch}{1.48}
\begin{tabular}{@{}P{1.6cm}P{1.6cm}P{2.0cm}P{2.4cm}@{}}
\toprule
$\delta_1$ & $\delta_2$ & \textbf{Anbieter} & \\
\midrule
$0{,}90$ & $0{,}90$ & $0{,}526$ & Erstzieher-Vorteil\\
$0{,}99$ & $0{,}99$ & $0{,}503$ & fast hälftig\\
$0{,}50$ & $0{,}50$ & $0{,}667$ & große Ungeduld\\
$0{,}90$ & $0{,}80$ & $\mathbf{0{,}714}$ & \textcolor{gruen}{geduldiger}\\
$0{,}80$ & $0{,}90$ & $0{,}357$ & \textcolor{akzent}{ungeduldiger}\\
\bottomrule
\end{tabular}}
\end{center}

\begin{center}
\begin{minipage}{0.9\textwidth}
\textbf{Geduld ist wertvoller als der erste Zug.} Bei $\delta_1 = 0{,}8$ gegen $\delta_2 = 0{,}9$ verliert der Anbieter trotz des Vortritts.
\end{minipage}
\end{center}

\schritt{Der Zusammenhang mit Satz 14.1}
Für $\delta_1 = \delta_2 = \delta \to 1$ strebt $x^{*} \to \tfrac12$.

\begin{center}
\begin{minipage}{0.88\textwidth}
\kursiv{Allgemeiner gilt: \textbf{Im Grenzfall verschwindender Ungeduld liefert das Rubinstein-Modell genau die Nash-Verhandlungslösung} --- mit den Abzinsungsfaktoren als Gewichten. Das ist bemerkenswert, weil beide Zugänge nichts miteinander zu tun haben: Der eine gibt Forderungen vor, der andere rechnet ein Gleichgewicht aus. Man nennt so etwas eine \Begriff{nichtkooperative Begründung} einer kooperativen Lösung.}
\end{minipage}
\end{center}

\schritt{Und ein Vorbehalt}
\begin{center}
\begin{minipage}{0.88\textwidth}
\kursiv{Das Gleichgewicht sagt vorher, dass das \kursiv{erste} Angebot angenommen wird --- es gibt keine Verhandlungsrunden. \textbf{Die beobachtete Wirklichkeit sieht anders aus}, und der Grund ist meist, dass die Beteiligten die Ungeduld des anderen nicht kennen. Modelle mit unvollständiger Information erzeugen Verzögerung und sind erheblich schwerer zu behandeln.}
\end{minipage}
\end{center}
\end{merksatz}

\section{Mehr als zwei}

\begin{definition}
\boxtitelm{Definition 14.3 --- Koalitionsspiel}
Eine \Begriff{charakteristische Funktion} ordnet jeder Teilmenge $S$ der Spieler den Wert $v(S)$ zu, den sie gemeinsam erreichen kann.
\end{definition}

\begin{merksatz}
\boxtitel{Der Leitfall mit drei Abteilungen}
Drei Abteilungen können ein Großgerät gemeinsam betreiben.

\begin{center}
{\small
\renewcommand{\arraystretch}{1.5}
\begin{tabular}{@{}P{2.2cm}P{1.6cm}P{2.2cm}P{1.6cm}@{}}
\toprule
\textbf{Gruppe} & $v$ & \textbf{Gruppe} & $v$\\
\midrule
$\{1\}$ & $0$ & $\{1,2\}$ & $90$\\
$\{2\}$ & $0$ & $\{1,3\}$ & $80$\\
$\{3\}$ & $0$ & $\{2,3\}$ & $70$\\
& & $\{1,2,3\}$ & $120$\\
\bottomrule
\end{tabular}}
\end{center}

\begin{center}
\begin{minipage}{0.9\textwidth}
\textbf{Allein erreicht niemand etwas. Je zwei erreichen viel. Alle drei erreichen am meisten.}
\end{minipage}
\end{center}
\end{merksatz}

\begin{definition}
\boxtitelm{Definition 14.4 --- der Kern}
Eine Aufteilung $\vek x$ liegt im \Begriff{Kern}, wenn
\[ \sum_i x_i = v(N) \qquad\text{und}\qquad \sum_{i\in S} x_i \ \ge\ v(S) \ \text{ für alle } S . \]
\end{definition}

\begin{merksatz}
\boxtitel{Der Kern im Leitfall}
\schritt{Die Bedingungen}
\[ x_1+x_2 \ge 90, \qquad x_1+x_3 \ge 80, \qquad x_2+x_3 \ge 70, \qquad \textstyle\sum x_i = 120 \]

\schritt{Die entscheidende Beobachtung}
Addiert man die drei Ungleichungen:
\[ 2\,(x_1+x_2+x_3) \ \ge\ 90+80+70 = 240 \]
\[ 2\cdot120 = 240 \]

\begin{center}
\begin{minipage}{0.9\textwidth}
\textbf{Gleichheit --- also muss jede einzelne Ungleichung mit Gleichheit gelten.}
\end{minipage}
\end{center}

\schritt{Die Auflösung}
\[ x_3 = 120-90 = 30, \qquad x_2 = 120-80 = 40, \qquad x_1 = 120-70 = 50 \]

\[ \text{Kern} = \{(50,\ 40,\ 30)\} \]

\kursiv{Probe:} $50+40 = 90$ \checkmark, $50+30 = 80$ \checkmark, $40+30 = 70$ \checkmark, Summe $120$ \checkmark
\end{merksatz}

\begin{warnung}
\boxtitelr{Der Kern kann leer sein}
Ändert man nur $v(\{1,2,3\})$ auf $80$ und setzt alle Paare auf $60$:
\[ 2\cdot80 = 160 \ <\ 180 = 60+60+60 \]

\begin{center}
\begin{minipage}{0.9\textwidth}
\textbf{Die Bedingungen sind unerfüllbar --- der Kern ist leer.}
\end{minipage}
\end{center}

\begin{center}
\begin{minipage}{0.88\textwidth}
\kursiv{In Worten: Je zwei Abteilungen können zusammen $60$ erwirtschaften, alle drei nur $80$. \textbf{Aus jeder Aufteilung kann sich ein Paar mit Gewinn absondern}, und das gilt für jede denkbare Aufteilung. Eine dauerhafte Einigung gibt es nicht; die Koalitionen wechseln fortwährend --- es ist das Gegenstück zum Condorcet-Zyklus aus Kapitel~12.}
\end{minipage}
\end{center}
\end{warnung}

\section{Der Shapley-Wert}

\begin{merksatz}
\boxtitel{Ein Wert, den es immer gibt}
Man stellt sich vor, die Spieler träten in zufälliger Reihenfolge hinzu. Jeder erhält, was er beim Hinzutreten beisteuert.

\begin{center}
\begin{minipage}{0.9\textwidth}
\[ \varphi_i = \frac{1}{n!}\sum_{\text{Reihenfolgen}} \bigl[v(\text{Vorgänger}+i) - v(\text{Vorgänger})\bigr] \]
\end{minipage}
\end{center}

\schritt{Satz 14.5 --- Shapley, 1953}
$\varphi$ ist die einzige Zuteilung mit den Eigenschaften:
\begin{center}
{\small
\renewcommand{\arraystretch}{1.5}
\begin{tabular}{@{}P{2.4cm}P{6.2cm}@{}}
\toprule
\textbf{Vollständig} & $\sum_i \varphi_i = v(N)$\\
\textbf{Symmetrie} & Wer denselben Beitrag leistet, erhält dasselbe.\\
\textbf{Nulltypen} & Wer nirgends etwas beisteuert, erhält nichts.\\
\textbf{Additiv} & Bei zwei zusammengelegten Spielen addieren sich die Werte.\\
\bottomrule
\end{tabular}}
\end{center}

\schritt{Die Rechnung im Leitfall}
\begin{center}
{\small
\renewcommand{\arraystretch}{1.46}
\begin{tabular}{@{}P{2.4cm}P{1.4cm}P{1.4cm}P{1.4cm}@{}}
\toprule
\textbf{Reihen\-folge} & $1$ & $2$ & $3$\\
\midrule
$1,2,3$ & $0$ & $90$ & $30$\\
$1,3,2$ & $0$ & $40$ & $80$\\
$2,1,3$ & $90$ & $0$ & $30$\\
$2,3,1$ & $50$ & $0$ & $70$\\
$3,1,2$ & $80$ & $40$ & $0$\\
$3,2,1$ & $50$ & $70$ & $0$\\
\midrule
\textbf{Summe} & $270$ & $240$ & $210$\\
\textbf{geteilt} & $\mathbf{45}$ & $\mathbf{40}$ & $\mathbf{35}$\\
\bottomrule
\end{tabular}}
\end{center}

\kursiv{Zur zweiten Zeile:} Bei der Reihenfolge $1,3,2$ steuert $3$ den Betrag $v(\{1,3\})-v(\{1\}) = 80$ bei und $2$ den Betrag $v(\{1,2,3\})-v(\{1,3\}) = 40$.

\kursiv{Probe:} $45+40+35 = 120$ \checkmark
\end{merksatz}

\begin{warnung}
\boxtitelr{Der Shapley-Wert liegt nicht immer im Kern}
\begin{center}
{\small
\renewcommand{\arraystretch}{1.5}
\begin{tabular}{@{}P{2.6cm}P{2.0cm}P{2.6cm}@{}}
\toprule
\textbf{Bedingung} & \textbf{Shapley} & \\
\midrule
$x_1+x_2 \ge 90$ & $85$ & \textcolor{akzent}{verletzt}\\
$x_1+x_3 \ge 80$ & $80$ & \checkmark\\
$x_2+x_3 \ge 70$ & $75$ & \checkmark\\
\bottomrule
\end{tabular}}
\end{center}

\begin{center}
\begin{minipage}{0.9\textwidth}
\textbf{Die Abteilungen $1$ und $2$ stünden sich besser, wenn sie $3$ ausschlössen} --- $90$ statt $85$.
\end{minipage}
\end{center}

\begin{center}
\begin{minipage}{0.88\textwidth}
\kursiv{Die beiden Begriffe antworten auf verschiedene Fragen. \textbf{Der Kern fragt, was haltbar ist; der Shapley-Wert, was angemessen ist.} Beide sind wohlbegründet und fallen hier auseinander: $(50,40,30)$ gegen $(45,40,35)$. Wer den Shapley-Wert als Verteilungsschlüssel festschreibt, muss damit rechnen, dass Teilgruppen ausscheren wollen.}
\end{minipage}
\end{center}

\schritt{Der Trost}
\begin{center}
\begin{minipage}{0.88\textwidth}
\kursiv{Für eine wichtige Klasse --- die \Begriff{konvexen} Spiele, bei denen der Beitrag eines Spielers mit der Größe der Gruppe wächst --- liegt der Shapley-Wert stets im Kern, und der Kern ist nie leer. \textbf{Viele Anwendungsfälle sind konvex}, etwa die Aufteilung von Fixkosten bei gemeinsamer Nutzung. Der Leitfall ist es nicht: Der Beitrag von $3$ zu $\{1,2\}$ beträgt $30$, zu $\{1\}$ aber $80$.}
\end{minipage}
\end{center}
\end{warnung}

\section{Der Bruchfall}

\begin{warnung}
\boxtitelr{Was die kooperative Theorie voraussetzt}
\begin{center}
{\small
\renewcommand{\arraystretch}{1.52}
\begin{tabular}{@{}P{2.4cm}P{6.2cm}@{}}
\toprule
\textbf{Bindende Verträge} & Absprachen müssen durchsetzbar sein. Ohne das gilt Teil~II.\\
\textbf{Übertragbarer Nutzen} & $v(S)$ als eine Zahl setzt voraus, dass sich Nutzen beliebig umverteilen lässt --- meist durch Geld.\\
\textbf{Keine Außenwirkung} & $v(S)$ hängt nicht davon ab, was die Übrigen tun. Bei Wettbewerb ist das falsch.\\
\bottomrule
\end{tabular}}
\end{center}

\schritt{Zur zweiten Zeile}
Satz 4.1 besagt, dass Nutzen nicht zwischen Personen vergleichbar ist.

\begin{center}
\begin{minipage}{0.88\textwidth}
\kursiv{Die charakteristische Funktion setzt genau das voraus --- sie gibt einer Gruppe \kursiv{einen} Wert, den sie aufteilen kann. \textbf{Das ist zulässig, wenn es ein Tauschmittel gibt}, dessen Nutzen für alle näherungsweise linear ist. Bei einer Budgetaufteilung ist das vertretbar, bei einer Verteilung von Diensten oder Organen nicht.}
\end{minipage}
\end{center}

\schritt{Zur dritten Zeile}
\begin{center}
\begin{minipage}{0.88\textwidth}
\kursiv{Schließen sich zwei Abteilungen zusammen, so verändert das häufig die Lage der dritten --- sie kann schlechter dastehen als zuvor. \textbf{Solche Spiele heißen Spiele in Partitionsfunktionsform}, und dort ist selbst der Kernbegriff mehrdeutig: Er hängt davon ab, was man der abgesonderten Gruppe über das Verhalten der Übrigen unterstellt. Das ist bis heute ungelöst.}
\end{minipage}
\end{center}
\end{warnung}

\section{Übungsaufgaben}

\begin{aufgabe}
\boxtitel{Aufgabe 14.1 --- Die Ausgangslage verbessern}
Das Klinikum kann seinen Konfliktwert durch einen Einzelvertrag von $20$ auf $40$ heben. Was bringt das, und was darf es kosten?

\loesung
\schritt{Schritt 1 --- Die neue Lösung}
\[ f(u_1) = (u_1-40)(70-u_1) \]
\[ f'(u_1) = (70-u_1)-(u_1-40) = 110-2u_1 = 0 \quad\Longrightarrow\quad u_1 = 55 \]

\[ \vek u^{*} = (55,\ 45) \]

\kursiv{Probe über den Überschuss:} $100-40-30 = 30$, hälftig $15$. Also $40+15 = 55$ und $30+15 = 45$. \checkmark

\schritt{Schritt 2 --- Der Gewinn}
\begin{center}
{\small
\renewcommand{\arraystretch}{1.5}
\begin{tabular}{@{}P{2.6cm}P{1.8cm}P{1.8cm}P{1.8cm}@{}}
\toprule
\textbf{Konfliktpunkt} & \textbf{K} & \textbf{M} & \textbf{Gewinn K}\\
\midrule
$(20,30)$ & $45$ & $55$ & --- \\
$(40,30)$ & $55$ & $45$ & $+10$\\
\bottomrule
\end{tabular}}
\end{center}

\begin{center}
\begin{minipage}{0.9\textwidth}
\textbf{Eine Verbesserung der Ausgangslage um $20$ bringt in der Verhandlung nur $10$} --- die andere Hälfte kommt dem Gegenüber zugute.
\end{minipage}
\end{center}

\schritt{Schritt 3 --- Die Regel}
\[ \frac{\partial u_1^{*}}{\partial d_1} = \frac12 \]

\begin{center}
\begin{minipage}{0.88\textwidth}
\kursiv{Bei linearer Grenze schlägt jede Verbesserung des Konfliktwerts zur Hälfte durch. \textbf{Der Einzelvertrag darf also höchstens $10$ kosten}, nicht $20$. Wer diese Halbierung übersieht, überschätzt den Wert von Rückfallmöglichkeiten regelmäßig um den Faktor zwei.}
\end{minipage}
\end{center}

\schritt{Schritt 4 --- Die Gegenmaßnahme des Marienhospitals}
Hebt es seinerseits von $30$ auf $50$:
\[ \text{Überschuss} = 100-40-50 = 10, \qquad \vek u^{*} = (45,\ 55) \]

\begin{center}
\begin{minipage}{0.9\textwidth}
\textbf{Beide haben aufgerüstet, und das Ergebnis ist dasselbe wie am Anfang --- nur wurde zweimal bezahlt.}
\end{minipage}
\end{center}

\begin{center}
\begin{minipage}{0.88\textwidth}
\kursiv{Die Verbesserung der Ausgangslage ist selbst ein Spiel, und es hat die Gestalt eines Gefangenendilemmas: \textbf{Beide investieren, und der Überschuss schrumpft von $50$ auf $10$.} Das ist der Kern dessen, was man in der Wirtschaftstheorie unbeschäftigte Kosten des Wettbewerbs nennt --- Aufwand, der keinen Wert schafft, sondern nur Verhandlungsmacht verschiebt.}
\end{minipage}
\end{center}
\end{aufgabe}

\begin{aufgabe}
\boxtitel{Aufgabe 14.2 --- Was Geduld wert ist}
Im Rubinstein-Modell mit $\delta_1 = 0{,}9$ und $\delta_2 = 0{,}8$: Wie viel erhält jeder, und wie ändert es sich, wenn Spieler 2 zuerst anbietet?

\loesung
\schritt{Schritt 1 --- Spieler 1 bietet zuerst}
\[ x^{*} = \frac{1-\delta_2}{1-\delta_1\delta_2} = \frac{1-0{,}8}{1-0{,}72} = \frac{0{,}2}{0{,}28} = \frac57 = 0{,}714 \]

\begin{center}
\begin{minipage}{0.9\textwidth}
\textbf{Spieler 1 erhält $71{,}4$ Prozent.}
\end{minipage}
\end{center}

\schritt{Schritt 2 --- Spieler 2 bietet zuerst}
Die Formel gilt entsprechend mit vertauschten Rollen:
\[ y^{*} = \frac{1-\delta_1}{1-\delta_1\delta_2} = \frac{0{,}1}{0{,}28} = \frac{5}{14} = 0{,}357 \]

Spieler 1 erhält dann $1-0{,}357 = 0{,}643$.

\schritt{Schritt 3 --- Der Vergleich}
\begin{center}
{\small
\renewcommand{\arraystretch}{1.5}
\begin{tabular}{@{}P{2.6cm}P{2.0cm}P{2.0cm}@{}}
\toprule
\textbf{Wer bietet} & \textbf{Spieler 1} & \textbf{Spieler 2}\\
\midrule
Spieler 1 & $0{,}714$ & $0{,}286$\\
Spieler 2 & $0{,}643$ & $0{,}357$\\
\midrule
\textbf{Unterschied} & $0{,}071$ & $0{,}071$\\
\bottomrule
\end{tabular}}
\end{center}

\begin{center}
\begin{minipage}{0.9\textwidth}
\textbf{Der Erstzieher-Vorteil beträgt sieben Prozentpunkte, der Geduldsvorteil des Spielers 1 aber über dreißig.}
\end{minipage}
\end{center}

\schritt{Schritt 4 --- Geduld gegen Vortritt}
Ist Spieler 2 der Geduldigere ($\delta_1=0{,}8$, $\delta_2=0{,}9$) und bietet Spieler 1 zuerst:
\[ x^{*} = \frac{0{,}1}{1-0{,}72} = 0{,}357 \]

\begin{center}
\begin{minipage}{0.9\textwidth}
\textbf{Spieler 1 erhält trotz des Vortritts nur $36$ Prozent.}
\end{minipage}
\end{center}

\begin{center}
\begin{minipage}{0.88\textwidth}
\kursiv{Der Unterschied in der Geduld beträgt nur ein Zehntel und kehrt das Ergebnis vollständig um. \textbf{Wer den Abschluss weniger dringend braucht, verhandelt besser} --- und das ist eine der wenigen Aussagen dieses Bandes, die sich unmittelbar in Handeln übersetzt: Man verbessere seine Rückfallmöglichkeit \kursiv{und} verlängere seinen Zeithorizont.}
\end{minipage}
\end{center}

\schritt{Schritt 5 --- Die Grenze}
\[ \lim_{\delta_1,\delta_2\to1} x^{*} \ \text{hängt vom Verhältnis der Ungeduld ab} \]

Mit $\delta_i = e^{-r_i \Delta}$ und $\Delta \to 0$:
\[ x^{*} \ \longrightarrow\ \frac{r_2}{r_1+r_2} \]

\begin{center}
\begin{minipage}{0.88\textwidth}
\kursiv{\textbf{Das ist die Nash-Verhandlungslösung mit den Ungeduldsraten als Gewichten.} Bei gleicher Ungeduld ergibt sich die hälftige Teilung. Der Erstzieher-Vorteil verschwindet, wenn die Runden beliebig kurz werden --- er war ein Artefakt der Zeittaktung.}
\end{minipage}
\end{center}
\end{aufgabe}

\begin{aufgabe}
\boxtitel{Aufgabe 14.3 --- Kern und Shapley vergleichen}
Drei Abteilungen teilen die Fixkosten eines Geräts. Allein kostet es jede $100$; je zwei zusammen $140$; alle drei $160$. Bestimmen Sie Kern und Shapley-Wert der \kursiv{Ersparnis}.

\loesung
\schritt{Schritt 1 --- Die charakteristische Funktion}
Die Ersparnis einer Gruppe ist die Differenz zwischen Einzelkosten und gemeinsamen Kosten.
\begin{center}
{\small
\renewcommand{\arraystretch}{1.48}
\begin{tabular}{@{}P{1.8cm}P{2.0cm}P{2.0cm}P{1.8cm}@{}}
\toprule
\textbf{Gruppe} & \textbf{einzeln} & \textbf{gemeinsam} & $v(S)$\\
\midrule
$\{i\}$ & $100$ & $100$ & $0$\\
$\{i,j\}$ & $200$ & $140$ & $60$\\
$\{1,2,3\}$ & $300$ & $160$ & $\mathbf{140}$\\
\bottomrule
\end{tabular}}
\end{center}

\schritt{Schritt 2 --- Der Kern}
\[ x_i+x_j \ \ge\ 60 \ \text{für alle Paare}, \qquad \textstyle\sum x_i = 140 \]

Addieren der drei Paarbedingungen:
\[ 2\cdot140 = 280 \ \ge\ 180 \quad \checkmark \]

\begin{center}
\begin{minipage}{0.9\textwidth}
\textbf{Der Kern ist nicht leer --- und diesmal ein ganzer Bereich.}
\end{minipage}
\end{center}

\kursiv{Die Bedingungen einzeln:} $x_i+x_j \ge 60$ bedeutet $x_k \le 80$. Zusammen mit $x_i \ge 0$:
\[ \text{Kern} = \Bigl\{\vek x \ : \ 0 \le x_i \le 80, \ \textstyle\sum x_i = 140\Bigr\} \]

\schritt{Schritt 3 --- Der Shapley-Wert}
Wegen der vollständigen Symmetrie:
\[ \varphi_i = \frac{140}{3} = 46{,}\overline{6} \]

\kursiv{Nachrechnung über eine Reihenfolge:} Der Erste steuert $0$ bei, der Zweite $60-0 = 60$, der Dritte $140-60 = 80$.
\[ \varphi_i = \frac{0+60+80}{3} = \frac{140}{3} \quad \checkmark \]

\begin{center}
\begin{minipage}{0.86\textwidth}
\kursiv{Die Abkürzung funktioniert hier, weil jeder Spieler in jeder Position gleich häufig vorkommt: \textbf{Bei Symmetrie genügt der Mittelwert der Beiträge über die Positionen}, statt alle $n!$ Reihenfolgen durchzugehen.}
\end{minipage}
\end{center}

\schritt{Schritt 4 --- Liegt er im Kern?}
\[ 46{,}\overline{6} \ \le\ 80 \quad \checkmark \]

\begin{center}
\begin{minipage}{0.9\textwidth}
\textbf{Ja.} Hier fallen beide Begriffe nicht auseinander.
\end{minipage}
\end{center}

\schritt{Schritt 5 --- Warum}
Das Spiel ist \kursiv{konvex}: Der Beitrag eines Spielers wächst mit der Gruppengröße.
\begin{center}
\renewcommand{\arraystretch}{1.48}
\begin{tabular}{@{}P{2.8cm}P{2.0cm}P{2.4cm}@{}}
\toprule
\textbf{tritt hinzu zu} & \textbf{Beitrag} & \\
\midrule
$\emptyset$ & $0$ & \\
einem anderen & $60$ & \\
zwei anderen & $80$ & wachsend\\
\bottomrule
\end{tabular}
\end{center}

\begin{center}
\begin{minipage}{0.88\textwidth}
\kursiv{Nach dem Satz von \Person{Shapley} (1971) liegt bei konvexen Spielen der Shapley-Wert stets im Kern, und der Kern ist nie leer. \textbf{Die Aufteilung gemeinsamer Fixkosten ist der Musterfall} --- und das ist der Grund, warum sich der Shapley-Wert gerade dort als Verteilungsschlüssel durchgesetzt hat, etwa bei Flughafengebühren nach Bahnlänge oder bei Gemeinschaftsleitungen.}
\end{minipage}
\end{center}

\schritt{Schritt 6 --- Die Kosten zurückgerechnet}
\[ \text{Kostenanteil}_i = 100 - 46{,}\overline{6} = 53{,}\overline{3} \]

\kursiv{Probe:} $3\cdot53{,}\overline{3} = 160$ \checkmark

\begin{center}
\begin{minipage}{0.88\textwidth}
\kursiv{Die gleichmäßige Aufteilung ergibt hier dasselbe, weil das Spiel symmetrisch ist. \textbf{Bei ungleichen Nutzungsanteilen weichen die Verfahren ab}, und dann ist der Shapley-Wert der begründete Schlüssel --- er verteilt nach dem durchschnittlichen Beitrag, nicht nach der Kopfzahl.}
\end{minipage}
\end{center}
\end{aufgabe}

\begin{geschichte}
\boxtitelg{Der andere Aufsatz von Nash}
\Person{John Nash} veröffentlichte 1950 zwei Arbeiten. Die eine begründete den Gleichgewichtsbegriff, die andere die Verhandlungslösung.

\begin{center}
\begin{minipage}{0.9\textwidth}
\kursiv{Die Verhandlungsarbeit entstand zuerst --- als Seminararbeit während seines Studiums in Princeton, noch bevor er von Neumann und Morgenstern gelesen hatte. \textbf{Er hatte den Aufsatz in einem Grundkurs über Außenwirtschaft geschrieben} und reichte ihn erst später zur Veröffentlichung ein.}
\end{minipage}
\end{center}

\schritt{Shapley}
\Person{Lloyd Shapley} (1923--2016) führte seinen Wert 1953 in seiner Dissertation ein.

\begin{center}
\begin{minipage}{0.9\textwidth}
\kursiv{Er war im Zweiten Weltkrieg als Soldat in China stationiert und erhielt dort eine Auszeichnung für die Entschlüsselung eines sowjetischen Wetterschlüssels. \textbf{Shapley arbeitete jahrzehntelang bei der RAND Corporation} und galt als einer der Wenigen, die den kooperativen Zweig weiterverfolgten, als fast alle zum nichtkooperativen übergingen.}
\end{minipage}
\end{center}

\schritt{Rubinstein}
\Person{Ariel Rubinstein} veröffentlichte sein Modell 1982 in einem Aufsatz von zwölf Seiten.

\begin{center}
\begin{minipage}{0.9\textwidth}
\kursiv{Es schloss eine Lücke, die seit Nash offen war: \textbf{Warum sollte das Ergebnis einer Verhandlung gerade die Nash-Lösung sein?} Rubinsteins Antwort --- weil sie das einzige teilspielperfekte Gleichgewicht eines naheliegenden Verhandlungsspiels ist --- gilt als Musterbeispiel dafür, wie ein kooperativer Begriff nichtkooperativ begründet wird.}
\end{minipage}
\end{center}

\schritt{Und eine Bemerkung des Verfassers}
Rubinstein hat sich wiederholt kritisch über die Anwendung der Spieltheorie geäußert.

\begin{center}
\begin{minipage}{0.9\textwidth}
\kursiv{Er hält sie für ein Werkzeug zum Nachdenken, nicht zum Vorhersagen, und warnt davor, aus Modellen Handlungsempfehlungen abzuleiten. \textbf{\enquote{Ich habe nie einen Fall gesehen, in dem die Spieltheorie jemandem geholfen hätte, eine bessere Entscheidung zu treffen}} --- das ist eine Äußerung von jemandem, der eines der meistzitierten Ergebnisse des Fachs bewiesen hat, und Kapitel~15 nimmt sie ernst.}
\end{minipage}
\end{center}
\end{geschichte}

\begin{forschung}
\boxtitelp{Kooperative Begriffe heute}
\begin{center}
{\small
\renewcommand{\arraystretch}{1.45}
\begin{tabular}{@{}P{2.8cm}P{5.4cm}@{}}
\toprule
\textbf{Kostenteilung} & Flughafengebühren, Gemeinschaftsleitungen, Rechenzentren.\\
\textbf{Erklärbarkeit} & Der Shapley-Wert als Maß für den Beitrag eines Merkmals zu einer Vorhersage.\\
\textbf{Netzwerke} & Wer in einem Verbund welchen Anteil an der Wertschöpfung hat.\\
\bottomrule
\end{tabular}}
\end{center}

\schritt{Zur zweiten Zeile}
Seit etwa 2017 wird der Shapley-Wert eingesetzt, um zu erklären, welchen Anteil ein einzelnes Merkmal an einer Modellvorhersage hat.

\begin{center}
\begin{minipage}{0.88\textwidth}
\kursiv{Die \enquote{Spieler} sind die Merkmale, $v(S)$ ist die Vorhersage bei Kenntnis nur der Merkmale in $S$. \textbf{Die vier Forderungen aus Satz 14.5 werden dabei zu Forderungen an eine Erklärung} --- Vollständigkeit, Gleichbehandlung, kein Beitrag ohne Wirkung, Zerlegbarkeit. Der Rechenaufwand ist allerdings exponentiell und muss genähert werden; und die Unterstellung, man könne Merkmale unabhängig weglassen, ist bei korrelierten Merkmalen falsch.}
\end{minipage}
\end{center}

\schritt{Zur ersten Zeile}
Das älteste Anwendungsbeispiel stammt von 1973: die Aufteilung der Kosten einer Landebahn auf Flugzeuge verschiedener Größe.

\begin{center}
\begin{minipage}{0.88\textwidth}
\kursiv{Der Shapley-Wert liefert dort eine geschlossene Formel: Jeder zahlt für das Stück Bahn, das er braucht, zu gleichen Teilen mit allen, die es ebenfalls brauchen. \textbf{Das Ergebnis ist so einleuchtend, dass es ohne Theorie gefunden worden wäre} --- und genau darin liegt seine Stärke als Prüfstein: Ein Verteilungsbegriff, der in klaren Fällen das Naheliegende liefert, ist in unklaren Fällen vertrauenswürdiger.}
\end{minipage}
\end{center}
\end{forschung}

\begin{merksatz}
\boxtitel{Was Sie jetzt können}
\begin{center}
\renewcommand{\arraystretch}{1.5}
\begin{tabular}{@{}P{3.0cm}P{5.6cm}@{}}
\toprule
\textbf{Satz 14.1} & Produkt der Zugewinne maximieren\\
\textbf{Satz 14.2} & Rubinstein: Geduld schlägt Vortritt\\
\textbf{Definition 14.4} & Kern: keine Gruppe schert aus\\
\textbf{Satz 14.5} & Shapley: mittlerer Beitrag\\
\textbf{Der Unterschied} & haltbar gegen angemessen\\
\bottomrule
\end{tabular}
\end{center}

\begin{center}
\begin{minipage}{0.88\textwidth}
\kursiv{Damit endet Teil~V. Das letzte Kapitel zieht Bilanz --- und nimmt die zweite Warnung des Vorworts ernst: dass Menschen nicht tun, was die Theorie sagt.}
\end{minipage}
\end{center}
\end{merksatz}
